\documentclass[12pt,a4paper]{article}
\usepackage{hcmut}

% numbers=none: bỏ gutter số dòng nên không cần lề trái dành cho số
\lstset{xleftmargin=0pt, framexleftmargin=0pt}

% Khối code breakable: cho phép ngắt trang trong listing để không bỏ trống cuối trang
\usepackage{tcolorbox}
\tcbuselibrary{listings,breakable,skins}
\newtcblisting{codeblock}[1][]{
  breakable, enhanced,
  colback=codeBg, colframe=codeFrame,
  boxrule=0.6pt, arc=1pt,
  left=4pt, right=3pt, top=2pt, bottom=2pt,
  before skip=4pt, after skip=4pt,
  listing only,
  listing options={style=academicListing, frame=none, backgroundcolor=\color{codeBg},
                   numbers=none, xleftmargin=0pt, framexleftmargin=0pt, #1}
}

% ==============================================================================
% CẤU HÌNH TIÊU ĐỀ ĐẦU TRANG & CHÂN TRANG (HEADER / FOOTER)
% ==============================================================================
% Thay đổi tên môn học (hiển thị góc trên bên phải trang nội dung)
\headercoursename{KIẾN TRÚC MÁY TÍNH (CO2008)}

% Thay đổi tiêu đề báo cáo (hiển thị góc dưới bên trái trang nội dung)
\footerreporttitle{Báo cáo Lab 2}

\begin{document}

% ==============================================================================
% 1. TRANG BÌA (KHÔNG HIỂN THỊ HEADER/FOOTER VÀ SỐ TRANG)
% ==============================================================================
\begin{titlepage}
\begin{center}

{\textbf{\large ĐẠI HỌC QUỐC GIA, THÀNH PHỐ HỒ CHÍ MINH}}\\[0.2cm]
{\textbf{\large TRƯỜNG ĐẠI HỌC BÁCH KHOA}}\\[0.2cm]
{\textbf{\large KHOA KHOA HỌC VÀ KỸ THUẬT MÁY TÍNH}}\\[1.0cm]

\begin{figure}[h!]
  \centering
  \includegraphics[width=3.6cm]{Images/bachkhoa-logo.png}
\end{figure}

\vspace{0.8cm}

% TÊN MÔN HỌC & MÃ MÔN
{\textbf{\Large KIẾN TRÚC MÁY TÍNH (CO2008)}}\\[0.4cm]

\noindent\rule{\linewidth}{0.6pt}\\[0.4cm]
% TÊN LOẠI BÁO CÁO & ĐỀ TÀI
{\textbf{\LARGE BÁO CÁO LAB 2}}\\[0.4cm]
{\textbf{\huge Branches and procedures}}\\[0.3cm]
\noindent\rule{\linewidth}{0.6pt}

\vspace{1.5cm}

% BẢNG THÔNG TIN GIẢNG VIÊN & SINH VIÊN
\begin{tabular}{rll}
  \Large\textbf{Sinh viên:} & \Large Hoàng Đức Tài & \Large - 2510219 \\
\end{tabular}

\vfill

% ĐỊA ĐIỂM VÀ THỜI GIAN
{\large TP. HỒ CHÍ MINH, THÁNG \the\month\ NĂM \the\year}

\end{center}
\end{titlepage}

\newpage

% ==============================================================================
% 2. MỤC LỤC (KHÔNG ĐÁNH SỐ TRANG)
% ==============================================================================
\pagestyle{empty}
\tableofcontents
\newpage

% ==============================================================================
% 3. NỘI DUNG CHÍNH CỦA BÁO CÁO (BẮT ĐẦU ĐÁNH SỐ TRANG 1 TỪ TRANG NỘI DUNG)
% ==============================================================================
\pagestyle{fancy}
\setcounter{page}{1}
\section{Exercise 1}

Mảng 5 phần tử khai báo bằng \texttt{.word 0:5}. \texttt{\$s0} giữ địa chỉ đầu mảng, \texttt{\$t0} là chỉ số phần tử đang xử lý:

\begin{codeblock}[language=MIPS]
  array: .word 0:5

  la $s0, array    # địa chỉ đầu mảng
  li $t0, 0        # chỉ số index
\end{codeblock}

Mỗi vòng lặp đọc một số nguyên, chuyển số đó về bội số của 4 gần nhất rồi lưu vào mảng.

Phần dư lấy bằng \texttt{andi \$t1, \$v0, 3}. Chia cho $4 = 2^2$ là bỏ đi hai bit thấp, nên phần dư nằm gọn trong hai bit đó; AND với $3 = 11_2$ giữ lại đúng hai bit ấy và xoá phần còn lại. Cách này chỉ cần một lệnh, không dùng đến \texttt{HI}/\texttt{LO} như \texttt{div}, và vì là bitwise nên số âm cũng ra phần dư $0$–$3$ như phép chia dư không âm.

Phần dư 0 thì giữ nguyên; dư 1 hoặc 2 thì trừ đi phần dư (làm tròn xuống); dư 3 thì cộng thêm 1 (làm tròn lên). Vì thế 35 thành 36, còn 34 và 33 thành 32.

\begin{codeblock}[language=MIPS]
process_loop:
  li $v0, 5                # syscall 5: đọc số nguyên
  syscall

  andi $t1, $v0, 3         # $t1 = n % 4
  sub $v0, $v0, $t1        # n -= n % 4: làm tròn xuống
  beq $t1, 3, round_up     # dư 3 thì bội số gần nhất nằm phía trên
  j next

round_up:
  addi $v0, $v0, 4         # n += 4: nhảy lên bội số kế tiếp

next:
  sll $t1, $t0, 2          # $t1 = 4i
  add $t1, $t1, $s0        # $t1 = base + 4i
  sw $v0, 0($t1)           # lưu vào phần tử thứ i
  addi $t0, $t0, 1         # i++
  bne $t0, 5, process_loop # đủ 5 phần tử thì dừng
\end{codeblock}

Vòng lặp thứ hai đọc lại mảng và in từng phần tử, mỗi số cách nhau một dấu cách:

\begin{codeblock}[language=MIPS]
  li $t0, 0         # i = 0: duyệt lại từ đầu mảng
print:
  li $v0, 1         # syscall 1: in số nguyên
  sll $a0, $t0, 2   # $a0 = 4i
  add $a0, $s0, $a0 # $a0 = base + 4i
  lw $a0, 0($a0)    # $a0 = phần tử thứ i
  syscall

  li $v0, 11        # syscall 11: in ký tự
  li $a0, ' '       # in dấu cách
  syscall

  addi $t0, $t0, 1  # i++
  bne $t0, 5, print
\end{codeblock}

Chương trình không in lời nhắc khi đọc, nên nhập lần lượt 35, 34, 33, 12, 7 thì in ra:

\begin{codeblock}[language={}]
36 32 32 12 8
\end{codeblock}

\section{Exercise 2}

Các chuỗi trong \texttt{.data} đảm nhiệm lời nhắc, thông báo lỗi và hai chuỗi kết quả:

\begin{codeblock}[language=MIPS]
prompt_a: .asciiz "a = "
prompt_b: .asciiz "b = "
err_msg: .asciiz "Input must be a positive integer. Please try again.\n"
gcd_output: .asciiz "GCD = "
lcm_output: .asciiz "\nLCM = "
\end{codeblock}

Mỗi số đọc vào được kiểm tra bằng \texttt{bgtz}: lớn hơn 0 thì sang nhãn \texttt{valid\_a} lưu vào thanh ghi, ngược lại in \texttt{err\_msg} rồi quay lại chính nhãn đọc để hỏi lại. Hai khối đọc giống nhau hoàn toàn, chỉ khác chuỗi lời nhắc và thanh ghi đích:

\begin{codeblock}[language=MIPS]
read_a:
  li $v0, 4         # syscall 4: in chuỗi
  la $a0, prompt_a  # in "a = "
  syscall

  li $v0, 5         # syscall 5: đọc số nguyên
  syscall
  bgtz $v0, valid_a # a > 0 thì hợp lệ

  li $v0, 4         # syscall 4: in chuỗi
  la $a0, err_msg   # thông báo nhập sai
  syscall

  j read_a          # hỏi lại từ đầu

valid_a:
  move $s0, $v0     # $s0 = a
\end{codeblock}

GCD tính bằng vòng lặp Euclid: mỗi bước thay cặp $(a, b)$ bằng $(b, a \bmod b)$, dừng khi $b = 0$ thì $a$ là GCD. Hai thanh ghi \texttt{\$t0} và \texttt{\$t1} đóng vai trò $a$ và $b$, phần dư lấy qua \texttt{div} + \texttt{mfhi}:

\begin{codeblock}[language=MIPS]
gcd:
  move $t0, $s0          # $t0 = a
  move $t1, $s1          # $t1 = b

gcd_loop:
  beqz $t1, print_output # b == 0: $t0 là GCD

  div $t0, $t1           # chia a cho b
  mfhi $t2               # $t2 = a % b

  move $t0, $t1          # a = b
  move $t1, $t2          # b = a % b

  j gcd_loop             # lặp lại với cặp mới
\end{codeblock}

LCM tính từ GCD theo $\text{lcm}(a, b) = a \times b / \gcd(a, b)$, ngay tại nhãn \texttt{print\_output} với \texttt{mul} rồi \texttt{div} + \texttt{mflo}:

\begin{codeblock}[language=MIPS]
print_output:
  mul $t1, $s0, $s1 # $t1 = a * b
  div $t1, $t0      # chia cho GCD
  mflo $t1          # $t1 = LCM
\end{codeblock}

Kết quả được in ra:

\begin{codeblock}[language=MIPS]
  li $v0, 4          # syscall 4: in chuỗi
  la $a0, gcd_output # in "GCD = "
  syscall

  li $v0, 1          # syscall 1: in số nguyên
  move $a0, $t0      # in GCD
  syscall

  li $v0, 4          # syscall 4: in chuỗi
  la $a0, lcm_output # in "\nLCM = "
  syscall

  li $v0, 1          # syscall 1: in số nguyên
  move $a0, $t1      # in LCM
  syscall
\end{codeblock}

Nhập $a = 6$, $b = 45$:

\begin{codeblock}[language={}]
a = 6
b = 45
GCD = 3
LCM = 90
\end{codeblock}

Nhập $a = -5$ (bị từ chối), sau đó $a = 6$, $b = 45$:

\begin{codeblock}[language={}]
a = -5
Input must be a positive integer. Please try again.
a = 6
b = 45
GCD = 3
LCM = 90
\end{codeblock}

\section{Exercise 3}

Mảng 10 phần tử và ba thanh ghi lưu trạng thái: \texttt{\$s1} là giá trị lớn nhất, \texttt{\$s2} là giá trị lớn thứ hai, \texttt{\$s3} là cờ báo đã có giá trị lớn thứ hai hay chưa. Phần tử đầu tiên được lưu vào mảng và cũng làm mốc khởi tạo cho \texttt{\$s1}:

\begin{codeblock}[language=MIPS]
  la $s0, arr               # $s0 = địa chỉ mảng
  li $s1, 0                 # $s1 = giá trị lớn nhất
  li $s2, 0                 # $s2 = giá trị lớn thứ hai
  li $s3, 0                 # $s3 = 1 khi đã có giá trị lớn thứ hai
  li $t0, 0                 # $t0 = chỉ số i

input_loop:
  beq $t0, 10, print_second # đủ 10 phần tử thì in kết quả

  li $v0, 4                 # syscall 4: in chuỗi
  la $a0, prompt1           # in "Please insert a element: "
  syscall

  li $v0, 5                 # syscall 5: đọc số nguyên
  syscall

  sll $t1, $t0, 2           # $t1 = 4i
  add $t1, $t1, $s0         # $t1 = base + 4i
  sw $v0, 0($t1)            # lưu vào phần tử thứ i

  beqz $t0, first_input     # phần tử đầu chưa có mốc so sánh
\end{codeblock}

Sau khi đã có phần tử đầu, mỗi số đọc vào rơi vào bốn trường hợp dưới đây. Số bằng \texttt{\$s1} bị bỏ qua vì nó cũng là giá trị lớn nhất, không phải giá trị lớn thứ hai:

\begin{codeblock}[language=MIPS]
  bgt $v0, $s1, new_max    # lớn hơn max: max cũ thành second
  beq $v0, $s1, next       # bằng max: bỏ qua
  beqz $s3, new_second     # chưa có second: nhận luôn
  bgt $v0, $s2, new_second # lớn hơn second: thay second
  j next

first_input:
  move $s1, $v0            # phần tử đầu làm max
  j next

new_max:
  move $s2, $s1            # second = max cũ
  move $s1, $v0            # max = phần tử mới
  li $s3, 1                # đã có second
  j next

new_second:
  move $s2, $v0            # second = phần tử mới
  li $s3, 1                # đã có second
  j next
\end{codeblock}

Chương trình in giá trị lớn thứ hai trước, sau đó là một vòng quét thứ hai tìm mọi chỉ số có phần tử bằng \texttt{\$s2}. Trong vòng quét đó, \texttt{\$t1} là cờ cho biết đã in chỉ số nào chưa: nó khởi tạo bằng 0, chỉ số khớp đầu tiên được in một mình rồi cờ bật lên 1; từ chỉ số thứ hai trở đi mới in thêm \texttt{", "} phía trước. Nhờ vậy kết quả không có dấu phẩy thừa ở đầu:

\begin{codeblock}[language=MIPS]
print_second:
  beqz $s3, no_second     # chưa có second: báo không tồn tại

  li $v0, 4               # syscall 4: in chuỗi
  la $a0, msg1            # chuỗi "Second largest value is "
  syscall

  li $v0, 1               # syscall 1: in số nguyên
  move $a0, $s2           # in giá trị lớn thứ hai
  syscall

  li $v0, 4               # syscall 4: in chuỗi
  la $a0, msg2            # chuỗi ", found in index "
  syscall

  li $t0, 0               # $t0 = chỉ số i
  li $t1, 0               # $t1 = 1 khi đã in ít nhất một chỉ số

find_second_indexes:
  beq $t0, 10, exit       # quét hết mảng thì kết thúc

  sll $t2, $t0, 2         # $t2 = 4i
  add $t2, $t2, $s0       # $t2 = base + 4i
  lw $t3, 0($t2)          # $t3 = phần tử thứ i

  bne $t3, $s2, find_next # không bằng second: bỏ qua
  beqz $t1, print_idx     # chỉ số đầu tiên: in không dấu phẩy

  li $v0, 4               # syscall 4: in chuỗi
  la $a0, msg3            # in ", " trước các chỉ số sau
  syscall

print_idx:
  li $v0, 1               # syscall 1: in số nguyên
  move $a0, $t0           # in chỉ số
  li $t1, 1               # đánh dấu đã in ít nhất một chỉ số
  syscall

find_next:
  addi $t0, $t0, 1        # i++
  j find_second_indexes   # tiếp tục quét
\end{codeblock}

Khi cờ \texttt{\$s3} vẫn là 0, nhãn \texttt{no\_second} in thông báo rồi thoát:

\begin{codeblock}[language=MIPS]
no_second:
  li $v0, 4    # syscall 4: in chuỗi
  la $a0, msg5 # "No second largest value exists."
  syscall
  j exit
\end{codeblock}

Nhập $1, 2, 7, -7, 3, 7, 4, 6, 8, 7$:

\begin{codeblock}[language={}]
Second largest value is 7, found in index 2, 5, 9.
\end{codeblock}

Với mảng toàn số bằng nhau, cờ \texttt{\$s3} không bao giờ được bật. Nhập $5, 5, 5, 5, 5, 5, 5, 5, 5, 5$:

\begin{codeblock}[language={}]
No second largest value exists.
\end{codeblock}

\section{Exercise 4}

Chương trình đọc một số nguyên, đưa giá trị sang \texttt{\$a0} làm tham số cho \texttt{factorial} và kiểm tra hai trường hợp trước khi gọi. Ngưỡng 12 là giới hạn của số nguyên 32 bit có dấu: $12! = 479\,001\,600$ vẫn nằm dưới $2^{31} - 1 = 2\,147\,483\,647$, còn $13! = 6\,227\,020\,800$ đã vượt:

\begin{codeblock}[language=MIPS]
input:
  li $v0, 5             # syscall 5: đọc số nguyên
  syscall

  move $a0, $v0         # $a0 = n: tham số cho factorial
  bltz $a0, error       # n < 0: không có giai thừa
  bgt $a0, 12, overflow # n > 12: kết quả vượt 32 bit
\end{codeblock}

Kết quả trả về trong \texttt{\$v0} được lưu sang \texttt{\$t0} rồi in ra:

\begin{codeblock}[language=MIPS]
  jal factorial
  move $t0, $v0 # $t0 = n!

  li $v0, 1     # syscall 1: in số nguyên
  move $a0, $t0 # in n!
  syscall

  j exit
\end{codeblock}

Giai thừa viết đệ quy theo $n! = n \times (n-1)!$, dừng ở $0! = 1$. Sau lời gọi \texttt{jal factorial}, \texttt{\$a0} đã bị tầng dưới sửa, nên giá trị $n$ của tầng hiện tại phải lưu vào stack cùng với \texttt{\$ra}. Khi tầng dưới trả về, \texttt{\$v0} đang giữ $(n-1)!$, nhân với $n$ vừa khôi phục để được kết quả:

\begin{codeblock}[language=MIPS]
factorial:
  beqz $a0, base_case # 0! = 1

  addi $sp, $sp, -4   # mở slot cho $ra
  sw $ra, 0($sp)      # lưu địa chỉ trả về
  addi $sp, $sp, -4   # mở slot cho n
  sw $a0, 0($sp)      # lưu n

  addi $a0, $a0, -1   # tham số mới: n - 1
  jal factorial       # $v0 = (n-1)!

  lw $t0, 0($sp)      # lấy lại n
  addi $sp, $sp, 4    # đóng slot của n
  lw $ra, 0($sp)      # khôi phục địa chỉ trả về
  addi $sp, $sp, 4    # đóng slot của $ra

  mul $v0, $v0, $t0   # n! = n * (n-1)!
  jr $ra

base_case:
  li $v0, 1           # 0! = 1
  jr $ra
\end{codeblock}

Hai nhãn lỗi chỉ in thông báo tương ứng rồi nhảy tới \texttt{exit}:

\begin{codeblock}[language=MIPS]
error:
  li $v0, 4    # syscall 4: in chuỗi
  la $a0, msg1 # "A factorial for your number can not be found."
  syscall
  j exit

overflow:
  li $v0, 4    # syscall 4: in chuỗi
  la $a0, msg2 # "Number is too large for 32-bit integer."
  syscall
  j exit
\end{codeblock}

Nhập 5, chương trình in:

\begin{codeblock}[language={}]
120
\end{codeblock}

Nhập $-2$, giá trị âm rẽ sang nhãn \texttt{error}:

\begin{codeblock}[language={}]
A factorial for your number can not be found.
\end{codeblock}

Nhập 13, giá trị vượt ngưỡng nên rẽ sang nhãn \texttt{overflow}:

\begin{codeblock}[language={}]
Number is too large for 32-bit integer.
\end{codeblock}


\end{document}
